% ═══════════════════════════════════════════════════════════════════════
\section{Travailler avec l'IA}
% ═══════════════════════════════════════════════════════════════════════

\begin{frame}{Le bon partage des rôles}
  \small
  \begin{columns}[T]
    \begin{column}{0.48\textwidth}
      \begin{Rem}{Ce que l'IA fait bien}
        \footnotesize
        \begin{itemize}
          \item Le squelette d'un document complet.
          \item La mise en forme : tableaux, listes, en-têtes.
          \item Les figures TikZ et les formules longues.
          \item Expliquer un message d'erreur.
          \item Proposer des variantes d'une question.
        \end{itemize}
      \end{Rem}
    \end{column}
    \begin{column}{0.48\textwidth}
      \begin{AuTableau}[Ce qui reste votre travail]
        \footnotesize
        \begin{itemize}
          \item Vérifier les mathématiques : calculs, corrigé, unités.
          \item Ajuster le niveau, le vocabulaire, le contexte québécois.
          \item Repérer l'anglais (Question, Points, Solution) et le
            point décimal.
          \item Refuser les extensions inutiles ou inventées.
          \item Comprendre chaque bloc avant de le garder.
        \end{itemize}
      \end{AuTableau}
    \end{column}
  \end{columns}
\end{frame}

\begin{frame}{L'anatomie d'un bon prompt}
  \small
  \begin{center}
  \begin{tabular}{lp{9.5cm}}
    \toprule
    Contexte & qui vous êtes, pour quel niveau, quelle situation \\
    Tâche & ce qu'il faut produire, et son contenu \\
    Contraintes techniques & classe, pdfLaTeX sur Overleaf, extensions
      permises, français, virgule décimale, longueur \\
    Format de sortie & un seul fichier complet et compilable, commenté \\
    Vérification & demander le corrigé, les étapes de calcul, la liste
      des extensions et leur rôle \\
    \bottomrule
  \end{tabular}
  \end{center}
  \vspace{4pt}
  Les crochets \champ{comme ceci} dans les prompts qui suivent sont à
  remplacer par vos propres choix.
\end{frame}

\begin{frame}{Prompt type : le squelette d'un document}
  \begin{Prompt}{Créer un document de départ}
    Tu es un enseignant de mathématiques \champ{au secondaire ou au
    collégial} au Québec et un expert de LaTeX.

    Écris un document LaTeX complet de classe article, qui compile sans
    erreur avec pdfLaTeX sur Overleaf. Il sera utilisé pour
    \champ{une fiche de notes sur les fonctions affines, 3\textsuperscript{e}
    secondaire} \champ{ou : un résumé sur les limites, Calcul différentiel au
    collégial}.

    Contraintes : texte en français ; extensions limitées à fontenc,
    babel (french), amsmath, amssymb, siunitx (virgule décimale), geometry
    et tikz ; aucune extension non standard ; pas de gras ni d'italique
    décoratifs.

    Ajoute un commentaire court à côté de chaque ligne du préambule pour
    expliquer son rôle. Donne uniquement le fichier .tex.
  \end{Prompt}
\end{frame}

\begin{frame}{Prompt type : corriger une erreur de compilation}
  \begin{Prompt}{Comprendre et réparer}
    Mon document LaTeX ne compile pas sur Overleaf (pdfLaTeX). Voici le
    message d'erreur complet :

    \champ{coller le message, par exemple : ! Missing \$ inserted.
    l.27 ...}

    Voici les lignes 20 à 30 de mon fichier :

    \champ{coller les lignes}

    Explique-moi en une ou deux phrases la cause de l'erreur, puis donne
    seulement les lignes corrigées. Ne réécris pas le reste du document.
  \end{Prompt}
  \small Donner le message exact et le passage concerné : sans eux,
  l'IA devine. Demander seulement les lignes corrigées : on garde le
  contrôle de son fichier.
\end{frame}

\begin{frame}{Prompts types : comprendre, modifier, dessiner}
  \begin{columns}[T]
    \begin{column}{0.49\textwidth}
      \begin{Prompt}{Expliquer un code}
        \scriptsize Voici un extrait de code LaTeX produit par une IA :
        \champ{extrait}. Explique ligne par ligne ce que fait chaque
        commande, en termes simples, pour une personne qui débute.
        Signale les lignes qui ne sont pas nécessaires.
      \end{Prompt}
      \vspace{2pt}
      \begin{Prompt}{Modification ciblée}
        \scriptsize Dans mon document, je veux \champ{ajouter une
        colonne « Points » au tableau de la section 2}. Donne seulement
        le bloc modifié, et indique où le placer.
      \end{Prompt}
    \end{column}
    \begin{column}{0.49\textwidth}
      \begin{Prompt}{Une figure TikZ}
        \scriptsize Écris le code TikZ d'une figure pour une fiche de
        \champ{3\textsuperscript{e} secondaire} : \champ{un triangle rectangle $ABC$,
        rectangle en $C$, avec $AC = 6$ cm et $BC = 8$ cm, le symbole
        d'angle droit et les mesures des côtés}. Uniquement
        l'environnement tikzpicture, avec des coordonnées entières et
        une échelle de \champ{0,5}, sans extension autre que tikz.
        Commente chaque ligne.
      \end{Prompt}
    \end{column}
  \end{columns}
\end{frame}

\begin{frame}{Avant de coller le code de l'IA}
  \begin{Rem}{Liste de vérification}
    \small
    \begin{enumerate}
      \item Je compile le code tel quel, avant toute modification.
      \item Je reconnais les trois parties : classe, préambule, corps.
      \item Chaque extension a un rôle que je peux nommer.
      \item Les nombres décimaux ont une virgule ; le texte est en
        français, y compris dans les en-têtes et les tableaux de points.
      \item J'ai refait moi-même les calculs du corrigé.
      \item Le niveau, le vocabulaire et le contexte conviennent à mes
        élèves ou à mes étudiants.
      \item Je sais où modifier un titre, une question, une formule, une
        coordonnée de figure.
    \end{enumerate}
  \end{Rem}
\end{frame}
